% ===========================================================================
% PART -1 -- every symbol and every object, from arithmetic upward.
% Nothing in this paper is used before it appears here. An audit of the first
% draft found thirty-one objects used without definition; all are below.
% ===========================================================================
\section{Every symbol, from arithmetic upward}
\label{sec:notation}

This section defines everything the paper uses. It assumes arithmetic --- adding,
multiplying, and the idea of a number --- and nothing else. A reader who is
fluent may skip to §\ref{sec:math}; a reader who is not should be able to finish
this section and then check every subsequent line.

\subsection{Numbers, inequalities, intervals}

$\R$ denotes the real numbers: everything on the number line, including fractions
and irrationals like $\sqrt 2$. $a<b$ means $a$ lies to the left of $b$.
$[a,b]$ is the set of all $x$ with $a \le x \le b$ (endpoints included);
$(a,b)$ excludes the endpoints. $|x|$ is the \emph{absolute value}: $x$ if
$x \ge 0$, and $-x$ otherwise. It is the distance from $x$ to $0$, so it is never
negative.

\step{One fact used repeatedly.} Multiplying an inequality by a \emph{positive}
number preserves it; multiplying by a \emph{negative} number \textbf{reverses}
it. $2<3$ but $-2>-3$. §\ref{sec:ci} depends on this at one specific step, and
the step is marked there.

\subsection{Functions}

A \emph{function} $f$ from a set $\mathcal{A}$ to a set $\mathcal{B}$, written
$f:\mathcal{A}\to\mathcal{B}$, assigns to each element of $\mathcal{A}$ exactly
one element of $\mathcal{B}$. $\mathcal{A}$ is the \emph{domain},
$\mathcal{B}$ the \emph{codomain}. ``Exactly one'' is the whole content: a rule
that sometimes gives two answers is not a function.

\begin{definition}[Monotone]
$f$ is \emph{strictly increasing} if $x<y$ implies $f(x)<f(y)$. Such a function
preserves order and nothing else: it can stretch and squeeze distances freely.
§\ref{sec:cardinal} uses exactly this property to test whether the paper's
central data are ordinal or not.
\end{definition}

\subsection{Sum and product notation}

\[
\sum_{i=1}^{n} x_i \;=\; x_1+x_2+\cdots+x_n,
\qquad
\prod_{i=1}^{n} x_i \;=\; x_1 \cdot x_2 \cdots x_n .
\]
The letter $i$ is a \emph{dummy}: $\sum_{i=1}^{3}x_i$ and $\sum_{j=1}^{3}x_j$ are
the same number. Three manipulation rules are used later and each is just
ordinary arithmetic written compactly.

\step{Rule 1 --- pulling out a constant.}
$\sum_i (c\,x_i) = c\sum_i x_i$, because $cx_1+cx_2+\cdots = c(x_1+x_2+\cdots)$.

\step{Rule 2 --- splitting a sum.}
$\sum_i (x_i+y_i) = \sum_i x_i + \sum_i y_i$, because addition can be reordered.

\step{Rule 3 --- reordering, and why it is safe here.} For a \emph{finite}
collection, adding the terms in any order gives the same total. Every sum in this
paper is finite, so grouping terms differently --- for instance grouping outcomes
by the value a function takes on them, as in the proof of \cref{lem:lin} --- never
changes the answer. (For infinite sums this is false in general; the paper never
needs one.)

\begin{definition}[Factorial]
$n! := 1\cdot 2\cdot 3 \cdots n$, with $0!:=1$ by convention. It counts the number
of ways to arrange $n$ distinct objects in a row: $n$ choices for the first, then
$n-1$, and so on. $3!=6$, and the six orderings are listed explicitly in
§\ref{sec:toy}.
\end{definition}

\subsection{Powers, roots, logarithms}

$x^2 := x\cdot x$; $\sqrt{x}$ is the non-negative number whose square is $x$;
$x^{1/2}=\sqrt x$. The \emph{logarithm} $\log x$ is the power to which the base
must be raised to give $x$. Only two properties are used:
$\log(ab)=\log a+\log b$ and $\log$ is strictly increasing. The second is why
correlating $\log$-transformed quantities in \cref{rem:normalise} is a statement
about \emph{proportional} rather than absolute co-movement.

\subsection{Max, min, argmax, sup, inf}

$\max_i x_i$ is the largest of the $x_i$; $\min_i x_i$ the smallest.

\begin{definition}[argmax]
$\argmax_{i} x_i$ is the \emph{index} at which the maximum occurs, not the value.
If $x=(4,9,2)$ then $\max = 9$ and $\argmax = 2$. When two entries tie, the
$\argmax$ is not unique; the paper's decision rules break ties by index order, and
exact ties in continuous scores occur with probability zero in the data.
\end{definition}

\begin{definition}[Supremum]
$\sup$ is the least upper bound: the smallest number that is $\ge$ every element
of a set. It coincides with $\max$ when the maximum is attained, and exists even
when it is not. $\sup\{1-\tfrac1n : n\ge1\}=1$ although $1$ is never attained.
\cref{thm:vetoinf} states $\sup_X \Delta(X)=\infty$, meaning: for any bound you
name, some candidate set exceeds it.
\end{definition}

\subsection{Sets, and the four operations used}

$\mathcal{A}\cup\mathcal{B}$ (union) contains everything in either;
$\mathcal{A}\cap\mathcal{B}$ (intersection) everything in both;
$\mathcal{A}\subseteq\mathcal{B}$ (subset) means every element of $\mathcal{A}$
is in $\mathcal{B}$; $\varnothing$ is the empty set; $|\mathcal{A}|$ is the count
of elements; $\mathcal{A}\setminus\mathcal{B}$ removes from $\mathcal{A}$
everything that is in $\mathcal{B}$.

\begin{lemma}[Union bound, the pointwise version]
\label{lem:union}
$\mathbf{1}\{x\in\bigcup_i \mathcal{A}_i\}\;\le\;\sum_i \mathbf{1}\{x\in\mathcal{A}_i\}$.
\end{lemma}
\begin{proof}
If $x$ is in none of the sets, both sides are $0$. If $x$ is in at least one, the
left side is $1$ and the right side is at least $1$ because it counts the sets
containing $x$. There is no third case.
\end{proof}
\Cref{thm:bonf} is this lemma plus taking expectations, and nothing else.

\subsection{Tuples and vectors}

An \emph{ordered tuple} $(x_1,\dots,x_k)$ is a list in which position matters:
$(1,2)\neq(2,1)$. The set of all such lists of $k$ real numbers is written $\R^k$,
and an element of it is a \emph{vector}.

Throughout this paper $k=4$ and the four positions are the four candidate
responses $A,B,C,D$. So a vector such as $(5.1,\,7.1,\,5.7,\,2.8)$ means: this
norm scores response $A$ at $5.1$, $B$ at $7.1$, and so on.

\begin{definition}[Addition and scaling]
$(x_1,\dots,x_k)+(y_1,\dots,y_k) := (x_1+y_1,\dots,x_k+y_k)$ and
$c\,(x_1,\dots,x_k) := (cx_1,\dots,cx_k)$. Both act position by position.
\end{definition}

\step{Why this is the right operation here.} A criterion with weight $w$ and
satisfaction profile $s=(s_A,s_B,s_C,s_D)$ contributes $w\,s$ to the score, and
several criteria contribute their sum. That is exactly \cref{def:score}, and it is
the only place vector arithmetic enters the construction.

\subsection{The dot product, the norm, and the angle}

\begin{definition}[Dot product]
$x\cdot y := \sum_{i=1}^{k} x_i y_i$, a single number.
\end{definition}

\begin{definition}[Norm]
$\|x\| := \sqrt{x\cdot x} = \sqrt{\sum_i x_i^2}$.
\end{definition}

\step{Why the norm is the length.} In two dimensions, $x=(x_1,x_2)$ is the corner
of a right triangle with legs $x_1$ and $x_2$; Pythagoras gives hypotenuse
$\sqrt{x_1^2+x_2^2}$, which is $\|x\|$. In $k$ dimensions the same formula is
taken as the \emph{definition} of length, and it behaves as length should:
$\|cx\|=|c|\,\|x\|$, and $\|x\|=0$ only when every coordinate is $0$.

\begin{theorem}[The dot product measures angle]
\label{thm:cosangle}
For non-zero $x,y$,
\[
\cos\theta \;=\; \frac{x\cdot y}{\|x\|\,\|y\|},
\]
where $\theta$ is the angle between $x$ and $y$.
\end{theorem}
\begin{proof}
\emph{Step 1.} The law of cosines for the triangle with sides $\|x\|$, $\|y\|$
and $\|x-y\|$ states
$\|x-y\|^2 = \|x\|^2+\|y\|^2-2\|x\|\|y\|\cos\theta$.

\emph{Step 2.} Expand the left side using the definition:
\[
\|x-y\|^2=\sum_i (x_i-y_i)^2=\sum_i x_i^2-2\sum_i x_iy_i+\sum_i y_i^2
=\|x\|^2-2(x\cdot y)+\|y\|^2 .
\]

\emph{Step 3.} Set the two expressions equal. The terms $\|x\|^2$ and $\|y\|^2$
cancel from both sides, leaving $-2(x\cdot y) = -2\|x\|\|y\|\cos\theta$.

\emph{Step 4.} Divide by $-2\|x\|\|y\|$.
\end{proof}

\begin{definition}[Cosine similarity]
$\cos(x,y) := \dfrac{x\cdot y}{\|x\|\|y\|}$, a number in $[-1,1]$.
$+1$ means same direction, $0$ perpendicular, $-1$ opposite.
\end{definition}

\begin{remark}[The single most important consequence for this paper]
Cosine similarity is \emph{blind to length}: $\cos(x,y)=\cos(2x,y)$ for any
scaling. That is precisely why it is the right tool once \cref{rem:normalise}
establishes that the length of $\Nn_i$ measures how many criteria a person
happened to rate rather than anything normative.
\end{remark}

\step{A worked cosine.} $x=(1,0)$, $y=(1,1)$. Then $x\cdot y = 1$,
$\|x\|=1$, $\|y\|=\sqrt2=1.4142$, so $\cos = 1/1.4142 = 0.7071$, which is
$\cos 45^\circ$. Correct: $(1,1)$ points at $45$ degrees to $(1,0)$.

\subsection{Centring, and what it removes}

\begin{definition}[Centred vector]
$\tilde x := x - \bar x\,\mathbf{1}$ where $\bar x = \tfrac1k\sum_i x_i$ and
$\mathbf{1}=(1,1,\dots,1)$. Equivalently, subtract the average from every entry.
\end{definition}

\step{Two facts, both immediate.} $\sum_i \tilde x_i = \sum_i x_i - k\bar x = 0$;
and if $y = x + c\mathbf{1}$ then $\tilde y = \tilde x$, because the shift moves
the average by exactly $c$ as well. \emph{Centring erases the common level and
keeps only the contrasts}, which is exactly what the decision uses.

\subsection{Matrices, rank, and the one matrix this paper needs}

A \emph{matrix} is a rectangular array of numbers; $M_{ij}$ is the entry in row
$i$, column $j$. A matrix acts on a vector by
$(Mx)_i := \sum_j M_{ij}x_j$: each output coordinate is a weighted sum of the
input coordinates.

\begin{definition}[Linear independence, span, rank]
Vectors $v_1,\dots,v_m$ are \emph{linearly independent} if the only way to make
$\sum_j c_j v_j = 0$ is to take every $c_j = 0$ --- none of them is redundant. Their
\emph{span} is the set of all $\sum_j c_j v_j$. The \emph{rank} of a matrix is the
number of linearly independent vectors among its outputs: the dimension of the
space it can reach.
\end{definition}

\begin{lemma}[The centring matrix and its rank]
\label{lem:centrank}
Let $C := I - \tfrac1k\mathbf{1}\mathbf{1}^\top$, so that $Cx = \tilde x$. Then
$C$ is a projection ($C^2=C$) and $\operatorname{rank}(C)=k-1$.
\end{lemma}
\begin{proof}
\emph{Step 1 --- $C$ does centre.} $(\mathbf{1}\mathbf{1}^\top x)_i = \sum_j x_j$
for every $i$, so $\tfrac1k \mathbf{1}\mathbf{1}^\top x = \bar x\mathbf{1}$, and
$Cx = x - \bar x \mathbf{1}$.

\emph{Step 2 --- it is idempotent.} Centring an already-centred vector changes
nothing, since its average is $0$. So $C(Cx)=Cx$, i.e. $C^2=C$.

\emph{Step 3 --- what it kills.} $C\mathbf{1} = \mathbf{1}-\mathbf{1}=0$, so the
one-dimensional space of constant vectors is sent to zero.

\emph{Step 4 --- what it keeps.} If $Cx=0$ then $x=\bar x\mathbf{1}$, i.e. $x$ is
constant. So the set sent to zero is \emph{exactly} the constants, of dimension
$1$.

\emph{Step 5 --- count.} The input space has dimension $k$; a one-dimensional
subspace is destroyed; the image therefore has dimension $k-1$. For $k=4$:
rank $3$.
\end{proof}

\Cref{lem:centrank} is the entire content of the first half of \cref{thm:probe}.

\subsection{Equivalence, quotients, and the sphere}

\begin{definition}[Equivalence relation]
A relation $\sim$ that is reflexive ($x\sim x$), symmetric ($x\sim y \Rightarrow
y \sim x$) and transitive ($x\sim y, y\sim z \Rightarrow x\sim z$). It partitions
a set into disjoint \emph{equivalence classes}: $[x] := \{y : y\sim x\}$.
\end{definition}

\step{The two used here.} (i) $u \sim u+c\mathbf{1}$: utilities differing by a
constant, which decide identically. (ii) $u \sim \lambda u$ for $\lambda>0$:
utilities differing by positive scale, which order identically.

\begin{definition}[Quotient]
The set of equivalence classes, written $\mathcal{A}/\!\sim$. Its dimension is
the dimension of $\mathcal{A}$ minus the dimension of what was collapsed.
\end{definition}

\begin{definition}[Sphere]
$S^{m} := \{x \in \R^{m+1} : \|x\|=1\}$, the set of unit-length vectors in
$\R^{m+1}$. It has dimension $m$: one constraint removes one degree of freedom.
$S^1$ is a circle in the plane; $S^2$ is the surface of a ball in three
dimensions.
\end{definition}

\step{Putting the two together, which is \cref{thm:probe}.} Starting from
$\R^4$: quotient by constants leaves dimension $3$; the unit-length
representatives of that $3$-dimensional space form $S^2$, of dimension $2$. Hence
``two free parameters''.

\subsection{Order statistics: median, quantile, percentile}

Sort the values as $x_{(1)}\le x_{(2)}\le\cdots\le x_{(n)}$.

\begin{definition}[Median]
The middle value: $x_{((n+1)/2)}$ if $n$ is odd, and the average of
$x_{(n/2)}$ and $x_{(n/2+1)}$ if $n$ is even. For $(2,4,4,4,6)$ it is $4$; for
$(2,4,4,6)$ it is $(4+4)/2=4$.
\end{definition}

\begin{definition}[Quantile]
The $q$-quantile is the value below which a fraction $q$ of the data lies. The
median is the $0.5$-quantile; $p_{25}$ and $p_{75}$ are the $0.25$- and
$0.75$-quantiles (\emph{quartiles}). ``$p_{95}=0.800$'' in §\ref{sec:est} means
$95\%$ of prompts have a relative margin below $0.800$.
\end{definition}

\begin{remark}[Why the paper reports quantiles beside means]
A mean can be moved arbitrarily by one extreme value; a median cannot. Reporting
both, as \cref{tab:vetoM} does with median $3.64$ and maximum $26.54$, exposes a
skew that a mean alone would hide --- and in that table the \emph{maximum} is the
load-bearing number, because \cref{thm:vetoinf} is about the worst case.
\end{remark}

\subsection{Rank correlation (Spearman)}

\begin{definition}[Spearman $\rho$]
Replace each value by its rank (smallest $=1$, next $=2$, \dots) in its own list,
then compute the ordinary correlation of the ranks.
\end{definition}

\step{Why the paper uses it for the gauge test.} It responds only to \emph{order},
so it answers ``do the two instruments rank the operators the same way'' without
being distracted by their different scales --- which is the exact question
§\ref{sec:gauge} asks. A Spearman of $+1.000$ means identical ordering; $0$ means
unrelated ordering.

\step{Worked.} Instrument 1 gives $(0.116, 0.051, 0.042, 0.023)$, ranks
$(4,3,2,1)$. Instrument 2 gives $(0.139,0.061,0.035,0.025)$, ranks $(4,3,2,1)$.
Identical ranks, so $\rho=+1.000$, even though no value matches.

\subsection{Bits, and what quantisation means}

\begin{definition}[Bit]
One \emph{bit} distinguishes two possibilities. $b$ bits distinguish $2^b$: one
bit gives $2$ levels, two bits $4$, three bits $8$, eight bits $256$.
\end{definition}

\begin{definition}[Quantisation to $b$ bits]
Given values in $[\ell,h]$, replace each $x$ by the nearest of the $2^b$ evenly
spaced levels $\ell + \tfrac{j}{2^b-1}(h-\ell)$, $j=0,\dots,2^b-1$.
\end{definition}

\step{Worked, at one bit.} With weights ranging over $[-2,8]$, one bit leaves the
two levels $\{-2, 8\}$: every negative weight becomes $-2$ and every positive one
$8$. Only the \emph{sign} survives. §\ref{sec:compile} reports that this
reproduces $89.8\%$ of the shipped decisions.

\subsection{Min--max normalisation}

\begin{definition}
$\hat x_i := \dfrac{x_i - \min_j x_j}{\max_j x_j - \min_j x_j} \in [0,1]$,
with $\hat x = 0$ at the worst entry and $1$ at the best.
\end{definition}

\step{Why §\ref{sec:agg} uses it, and what it assumes.} Comparing utility across
people requires putting them on a common scale, and nothing in the elicitation
supplies one. Min--max makes every person's best option worth $1$ and worst worth
$0$ \emph{by construction}, which is a stipulation, not a measurement. It is
stated as such: the Blackwell deficiency of $0.1298$ is a number relative to that
normalisation, and a different interpersonal scale would give a different figure.

\subsection{Conditional probability and conditional expectation}

\begin{definition}[Conditional probability]
$\Prob(E\mid F) := \dfrac{\Prob(E\cap F)}{\Prob(F)}$ when $\Prob(F)>0$:
restrict attention to the outcomes in $F$ and renormalise so they sum to $1$.
\end{definition}

\begin{definition}[Conditional expectation]
$\E[X \mid F] := \dfrac{\sum_{\omega\in F} X(\omega)\Prob(\omega)}{\Prob(F)}$,
the average of $X$ over $F$ alone.
\end{definition}

\begin{lemma}[Law of total expectation, the two-event case]
\label{lem:tower}
For an event $F$ with complement $F^c$,
\[
\E[X] \;=\; \E[X\mid F]\,\Prob(F) \;+\; \E[X\mid F^c]\,\Prob(F^c).
\]
\end{lemma}
\begin{proof}
Split the defining sum by whether the outcome lies in $F$:
\[
\E[X]=\sum_{\omega\in F}X(\omega)\Prob(\omega)+\sum_{\omega\in F^c}X(\omega)\Prob(\omega).
\]
Multiply and divide the first sum by $\Prob(F)$ and the second by $\Prob(F^c)$,
then apply the definition above to each. (If $\Prob(F)=0$ the term is absent and
the identity is trivial.)
\end{proof}

\Cref{lem:tower} with $F=\{\text{the decision flipped}\}$ and $X = \Reg$ is
exactly \cref{lem:coarsen}, since $\Reg=0$ on $F^c$ so the second term vanishes.

\subsection{Independent and identically distributed}

\emph{Identically distributed} means every observation is drawn from the same
$\Prob$. \emph{Independent} means no observation carries information about
another. Written \emph{iid}. §\ref{sec:cluster} exists because in this data the
observations are identically distributed but \textbf{not} independent, and only
the second half fails.

\subsection{Markov kernels}

\begin{definition}[Markov kernel]
A rule $K$ that, given an input $z$, produces a \emph{probability distribution}
over outputs, written $K(\cdot\mid z)$. A deterministic function is the special
case that puts all the probability on one output, written $\delta_{f(z)}$.
\end{definition}

\step{Why the paper needs it.} \Cref{thm:blackwell-proof} says a garbling cannot
help. Aggregation is deterministic, hence a garbling with the degenerate kernel
$\delta_{\Agg(z)}$; the theorem then applies verbatim without any probabilistic
aggregation being involved.

\subsection{Integrals, only as much as $\Phi$ requires}

The paper uses the normal distribution function
$\Phi(z)=\int_{-\infty}^{z}\tfrac{1}{\sqrt{2\pi}}e^{-t^2/2}\,dt$. The integral
sign means \emph{area under the curve to the left of $z$}. The curve
$\tfrac{1}{\sqrt{2\pi}}e^{-t^2/2}$ is the familiar bell; its total area is $1$, so
$\Phi(z)$ is the fraction of the bell lying left of $z$.

\begin{definition}[Quantile of the normal]
$z_q$ is the point with $\Phi(z_q)=q$. The three used are
$z_{0.975}=1.9600$, $z_{0.80}=0.8416$, $z_{0.99}=2.3263$. By the symmetry of the
bell about $0$, $z_{1-q}=-z_q$; this symmetry is the step used in
\cref{thm:mde}.
\end{definition}

\begin{remark}[No calculus is needed beyond this]
The paper never differentiates or integrates anything. $\Phi$ enters only as a
table of three numbers, and every use of it is an instance of ``how far out in
the bell is this''.
\end{remark}

\subsection{Sigmoid and logit, for the instrument}

The judge produces a number $\lambda$ (a difference of log-probabilities) and it
is mapped to $(0,1)$ by the \emph{sigmoid}
$\sigma(\lambda)=1/(1+e^{-\lambda})$, which is strictly increasing with
$\sigma(0)=0.5$.

\mesoboundary{$\sigma(\lambda)$ lies in $(0,1)$ and is therefore easily mistaken
for a probability. It is not one. It is a monotone re-expression of a model's
internal score, calibrated against nothing. Every use of $\sat$ in this paper
carries that caveat, and §\ref{sec:instrument} measures its consequence rather
than assuming it away.}

\subsection{A table of every symbol}

\begin{table}[H]\centering\footnotesize
\begin{tabular}{llll}
\toprule
symbol & meaning & first defined & \\
\midrule
$\mathbf{1}\{P\}$ & indicator, $1$ if $P$ true & §\ref{sec:math} &\\
$\sum,\ \prod,\ n!$ & sum, product, factorial & §\ref{sec:notation} &\\
$\argmax$ & the index of the maximum & §\ref{sec:notation} &\\
$\sup,\ \inf$ & least upper / greatest lower bound & §\ref{sec:notation} &\\
$x\cdot y,\ \|x\|$ & dot product, length & §\ref{sec:notation} &\\
$\cos(x,y)$ & cosine similarity & §\ref{sec:notation} &\\
$\tilde x$ & centred vector & §\ref{sec:notation} &\\
$C$ & centring matrix, rank $k-1$ & \cref{lem:centrank} &\\
$S^{m}$ & unit sphere, dimension $m$ & §\ref{sec:notation} &\\
$\Prob,\ \E,\ \Var,\ \Cov,\ \Corr$ & probability, mean, variance, covariance, correlation & §\ref{sec:math} &\\
$\se$ & standard error $=\operatorname{sd}/\sqrt n$ & §\ref{sec:se} &\\
$s^2$ & sample variance, divisor $n-1$ & \cref{lem:bessel} &\\
$\Phi,\ z_q$ & normal cdf and its quantiles & §\ref{sec:notation} &\\
$\rho$ & intraclass correlation \emph{or} Spearman & §\ref{sec:cluster}, §\ref{sec:notation} &\\
$\mathrm{DEFF}$ & design effect $1+(R-1)\rho$ & \cref{thm:deff} &\\
$\mathrm{MDE}$ & minimum detectable effect & \cref{thm:mde} &\\
$\mathrm{rel}$ & reliability & §\ref{sec:atten} &\\
$\resp,\ X,\ k$ & response space, probe set, its size ($=4$) & §\ref{sec:object} &\\
$\sat,\ \wt,\ V,\ \Pi$ & satisfaction, weight, veto set, ranking & \cref{def:tensors} &\\
$Y_p,\ \Nn_i$ & collective and individual induced scores & \cref{def:score}, \cref{def:ni} &\\
$\Agg,\ G$ & aggregation rule and its output & \cref{eq:chain} &\\
$\Reg$ & normalised regret & \cref{def:regret} &\\
$\pi$ & flip rate & §\ref{sec:est} &\\
$m$ & decision margin & \cref{prop:cap} &\\
$\Delta$ & the penalty a veto requires & \cref{thm:veto} &\\
$\phi_i$ & Shapley influence & §\ref{sec:c1} &\\
$v(S)$ & characteristic function of the coalition game & §\ref{sec:c1} &\\
$K(\cdot\mid z)$ & Markov kernel & §\ref{sec:notation} &\\
$\sigma(\lambda)$ & sigmoid & §\ref{sec:notation} &\\
\bottomrule
\end{tabular}
\caption{Every symbol used in the paper. $\rho$ is overloaded by convention; the
two uses never appear in the same expression and are disambiguated at each site.}
\label{tab:symbols}
\end{table}
